\documentclass[aspectratio=169]{beamer}
\usepackage{hyperref}
\usepackage[utf8]{inputenc}
\usepackage[T1]{fontenc}
\usepackage[spanish]{babel}
\usepackage{csquotes}
\usepackage{tribhuvan}

% other packages
\usepackage{latexsym,amsmath,xcolor,multicol,booktabs,calligra}
\usepackage{graphicx,listings,textcomp}
\usepackage{tikz}
\usetikzlibrary{arrows.meta,positioning,shapes.geometric}
% Resaltado de sintaxis para listados de ficheros (pixi.toml, recipe.yaml...),
% precomputado con `pygmentize -f latex` para no depender de -shell-escape.
\usepackage{fvextra}
\makeatletter
\def\PY@reset{\let\PY@it=\relax \let\PY@bf=\relax
  \let\PY@ul=\relax \let\PY@tc=\relax
  \let\PY@bc=\relax \let\PY@ff=\relax}
\def\PY@tok#1{\csname PY@tok@#1\endcsname}
\def\PY@toks#1+{\ifx\relax#1\empty\else
  \PY@tok{#1}\expandafter\PY@toks\fi}
\def\PY@do#1{\PY@bc{\PY@tc{\PY@ul{\PY@it{\PY@bf{\PY@ff{#1}}}}}}}
\def\PY#1#2{\PY@reset\PY@toks#1+\relax+\PY@do{#2}}

\def\PY@color#1{\def\PY@tc##1{\textcolor[rgb]{#1}{##1}}}
\@namedef{PY@tok@w}{\PY@color{0.73,0.73,0.73}}
\@namedef{PY@tok@c}{\let\PY@it=\textit\PY@color{0.24,0.48,0.48}}
\@namedef{PY@tok@c1}{\let\PY@it=\textit\PY@color{0.24,0.48,0.48}}
\@namedef{PY@tok@c+c1}{\let\PY@it=\textit\PY@color{0.24,0.48,0.48}}
\@namedef{PY@tok@k}{\let\PY@bf=\textbf\PY@color{0.00,0.50,0.00}}
\@namedef{PY@tok@k+kc}{\let\PY@bf=\textbf\PY@color{0.00,0.50,0.00}}
\@namedef{PY@tok@o}{\PY@color{0.40,0.40,0.40}}
\@namedef{PY@tok@p}{\PY@color{0.20,0.20,0.20}}
\@namedef{PY@tok@n}{\PY@color{0.10,0.09,0.49}}
\@namedef{PY@tok@n+nt}{\let\PY@bf=\textbf\PY@color{0.00,0.50,0.00}}
\@namedef{PY@tok@n+nv}{\PY@color{0.10,0.09,0.49}}
\@namedef{PY@tok@l}{\PY@color{0.73,0.13,0.13}}
\@namedef{PY@tok@l+s+s2}{\PY@color{0.73,0.13,0.13}}
\@namedef{PY@tok@l+lScalar+lScalarPlain}{\PY@color{0.73,0.13,0.13}}
\@namedef{PY@tok@p+pIndicator}{\PY@color{0.40,0.40,0.40}}
\@namedef{PY@tok@err}{\PY@color{0.89,0.00,0.00}}
\@namedef{PY@tok@s}{\PY@color{0.73,0.13,0.13}}
\@namedef{PY@tok@s+se}{\PY@color{0.73,0.13,0.13}}
\@namedef{PY@tok@m}{\PY@color{0.40,0.40,0.40}}
\@namedef{PY@tok@nt}{\let\PY@bf=\textbf\PY@color{0.00,0.50,0.00}}

\def\PYZbs{\char`\\}\def\PYZus{\char`\_}
\def\PYZob{\char`\{}\def\PYZcb{\char`\}}
\def\PYZca{\char`\^}\def\PYZam{\char`\&}
\def\PYZlt{\char`\<}\def\PYZgt{\char`\>}
\def\PYZsh{\char`\#}\def\PYZpc{\char`\%}
\def\PYZdl{\char`\$}\def\PYZhy{\char`\-}
\def\PYZsq{\char`\'}\def\PYZdq{\char`\"}
\def\PYZti{\char`\~}\def\PYZat{@}
\def\PYZlb{[}\def\PYZrb{]}
\makeatother


\author[Francisco Martín Rico, Esteve Fernández, Juan Sebastián Cely, Francisco Miguel Moreno]{%
\begin{tabular}{cc}
Francisco Martín Rico & Esteve Fernández \\
\texttt{\footnotesize fmrico@gmail.com} & \texttt{\footnotesize esteve.fernandez@gmail.com} \\[1em]
Juan Sebastián Cely & Francisco Miguel Moreno \\
\texttt{\footnotesize juan.cely@urjc.es} & \texttt{\footnotesize franciscom.moreno@urjc.es}
\end{tabular}%
}
\title[Pixi, RoboStack y Prefix.dev]{Desde que uso Pixi estoy tan Pichi}
\subtitle{Entornos reproducibles para ROS 2 con Pixi, RoboStack y Prefix.dev}
\institute{}

\date{}
\setbeamertemplate{date}{}
\setbeamertemplate{navigation symbols}{} % ocultar iconos de navegación

% defs
\definecolor{deepblue}{rgb}{0,0,0.5}
\definecolor{deepred}{rgb}{0.6,0,0}
\definecolor{deepgreen}{rgb}{0,0.5,0}
\definecolor{halfgray}{gray}{0.55}

\lstset{
    basicstyle=\ttfamily\small,
    keywordstyle=\bfseries\color{deepblue},
    emphstyle=\ttfamily\color{deepred},
    stringstyle=\color{deepgreen},
    numbers=none,
    breaklines=true,
    showstringspaces=false,
    rulesepcolor=\color{red!20!green!20!blue!20},
    frame=shadowbox,
    literate=
      {á}{{\'a}}1 {é}{{\'e}}1 {í}{{\'i}}1 {ó}{{\'o}}1 {ú}{{\'u}}1
      {Á}{{\'A}}1 {É}{{\'E}}1 {Í}{{\'I}}1 {Ó}{{\'O}}1 {Ú}{{\'U}}1
      {ñ}{{\~n}}1 {Ñ}{{\~N}}1
      {ü}{{\"u}}1 {Ü}{{\"U}}1
      {¿}{{\textquestiondown}}1 {¡}{{\textexclamdown}}1
}

% Frame auxiliar para preguntas/mini-retos interactivos al público
\newcommand{\preguntaframe}[1]{%
\begin{frame}
  \begin{center}
  \begin{beamercolorbox}[sep=1em,center]{block title}
  \Large \textbf{?} \; #1
  \end{beamercolorbox}
  \end{center}
\end{frame}%
}

% Frame auxiliar para marcar el inicio de un bloque práctico
\newcommand{\practicaframe}[2]{%
\begin{frame}
  \centering
  \vfill
  {\Large \textbf{PRÁCTICA}}\\[1em]
  {\large #1}\\[0.5em]
  {\color{halfgray} #2}
  \vfill
\end{frame}%
}

\begin{document}

%-------------------------------------------------
\begin{frame}
    \titlepage
\end{frame}

\begin{frame}[shrink=20]{Agenda}
    \tableofcontents[sectionstyle=show,subsectionstyle=show/shaded/hide,subsubsectionstyle=show/shaded/hide]
\end{frame}

%=================================================
\section[Introducción]{Introducción: ¿por qué ROS~2 necesita algo mejor?}
%=================================================

\begin{frame}{Antes de empezar\ldots}
\begin{itemize}
    \item ¿Cuántos habéis tenido un \lstinline|rosdep install| que falla distinto en vuestro portátil que en el del compañero?
    \item ¿Cuántos tenéis un \texttt{Dockerfile} de ROS~2 de más de 150 líneas solo para fijar versiones?
    \item ¿Cuántos tenéis instalado más de un \lstinline|/opt/ros/<distro>| en la misma máquina a la vez?
\end{itemize}
\vfill
\begin{block}{}
\centering No sois los únicos. Este workshop va de eso.
\end{block}
\end{frame}

\begin{frame}{El problema, en tres capas}
\begin{block}{Dependencias}
ROS~2 vía \texttt{apt} está atado a una versión de Ubuntu concreta. Mezclar ROS con librerías de ML modernas (PyTorch, CUDA reciente) choca con lo que trae el repositorio del sistema.
\end{block}
\begin{block}{Reproducibilidad}
\lstinline|apt install ros-jazzy-desktop| hoy y dentro de 6 meses puede instalar builds distintos. No hay lockfile. \texttt{rosdep} resuelve llamando a gestores distintos (\texttt{apt}, \texttt{pip}\ldots) sin fijar versiones.
\end{block}
\begin{block}{Overlays}
El modelo underlay/overlay de \texttt{colcon} (\lstinline|source /opt/ros/.../setup.bash| $\rightarrow$ \lstinline|source install/setup.bash|) es manual y frágil.
\end{block}
\end{frame}

\begin{frame}{¿Y Docker no resolvía esto?}
\begin{itemize}
    \item Docker aísla, pero \textbf{no fija automáticamente el contenido}: un \lstinline|apt-get install| dentro del Dockerfile sigue sin estar reproducido a menos que se congele todo a mano.
    \item Overhead de build/caché por capas $\rightarrow$ iteración lenta en desarrollo.
    \item No es multiplataforma ``gratis'': desarrollar en macOS/Windows para un target Linux añade capas de virtualización.
\end{itemize}
\vfill
\begin{alertblock}{Mensaje clave}
Docker sigue siendo válido para el despliegue final, pero no debería ser la \textbf{única} fuente de verdad de las dependencias del proyecto.
\end{alertblock}
\end{frame}

\begin{frame}{La idea central de hoy}
\begin{center}
\Large
¿Y si tuviéramos el equivalente a un \lstinline|package-lock.json|\\ o un \lstinline|Cargo.lock|\\[0.3em]
\textbf{para ROS~2}?
\vspace{1em}

\normalsize
Que funcione igual en tu portátil, en el del compañero y en CI.\\
Sin \texttt{sudo}. Sin depender de la versión de Ubuntu.\\
Sin tener que aprender Docker desde cero.
\end{center}
\end{frame}

\begin{frame}{Las tres piezas}
\begin{columns}[T]
\column{0.32\linewidth}
\begin{block}{Pixi}
Gestor de entornos reproducible: lockfile, tasks, multi-entorno.
\end{block}
\column{0.32\linewidth}
\begin{block}{RoboStack}
ROS~2 empaquetado como paquetes conda, multiplataforma.
\end{block}
\column{0.32\linewidth}
\begin{block}{Prefix.dev}
Dónde vive y se publica todo: canales conda, públicos y privados.
\end{block}
\end{columns}
\vfill
\centering Hoy vamos a recorrer las tres, de teoría a práctica.
\end{frame}

%=================================================
\section[¿Qué es Pixi?]{¿Qué es Pixi?}
%=================================================

\begin{frame}{¿Qué es Pixi?}
\begin{itemize}
    \item Gestor de entornos y paquetes \textbf{multiplataforma y multi-lenguaje}, escrito en Rust.
    \item Construido sobre \textbf{Rattler}: librerías Rust que reimplementan el ecosistema conda (resolución, instalación, canales). Rattler no es una CLI, es el motor; Pixi es la interfaz de usuario.
    \item Creado por el equipo de \textbf{prefix.dev} (los mismos autores de \textbf{mamba}).
\end{itemize}
\vfill
\begin{alertblock}{No es ``conda pero más rápido''}
Repiensa el flujo de trabajo: sin entorno base global, workspace-centric, lockfile nativo de primera clase.
\end{alertblock}
\end{frame}

\preguntaframe{¿Quién ha usado conda o mamba antes? ¿Qué es lo que más odiáis?}

\begin{frame}{Pixi frente a conda: diferencias de fondo}
\begin{itemize}
    \item Sin ``entorno base'' tipo Miniconda/Anaconda: Pixi es un binario estático, nada que activar a nivel de sistema.
    \item \textbf{Workspace-centric}: todo vive en una carpeta con \lstinline|pixi.toml| + \lstinline|pixi.lock|, autocontenida, clonable, versionable en git.
    \item Lockfile nativo de primera clase (\lstinline|pixi.lock|) — algo que conda nunca tuvo.
    \item Combina en el mismo manifiesto dependencias \textbf{conda} (C++, CUDA, binarios de sistema) y \textbf{PyPI} (resueltas con un motor basado en \texttt{uv}).
\end{itemize}
\end{frame}

\begin{frame}{Pixi vs. el resto del ecosistema}
\centering
\begin{tabular}{lcccc}
\toprule
\textbf{Feature} & \textbf{Pixi} & \textbf{Conda} & \textbf{Pip} & \textbf{Poetry} \\
\midrule
Lockfile reproducible        & \checkmark & -- & -- & \checkmark \\
Task runner integrado        & \checkmark & -- & -- & -- \\
Multi-lenguaje (no solo Py)  & \checkmark & \checkmark & -- & -- \\
Multi-entorno en un proyecto & \checkmark & parcial & -- & parcial \\
\bottomrule
\end{tabular}
\vfill
\begin{center}
\textit{Pixi resuelve e instala un entorno desde cero $\sim 3\times$ más rápido que micromamba y $>$10$\times$ más rápido que conda clásico.}
\end{center}
\end{frame}

\begin{frame}[fragile]{Estructura de un \texttt{pixi.toml}}
\begin{filelisting}
\begin{Verbatim}[commandchars=\\\{\}]
\PY{k}{[}\PY{k}{workspace}\PY{k}{]}
\PY{n}{name}\PY{+w}{ }\PY{o}{=}\PY{+w}{ }\PY{l+s+s2}{\PYZdq{}}\PY{l+s+s2}{my\PYZus{}ros2\PYZus{}project}\PY{l+s+s2}{\PYZdq{}}
\PY{n}{version}\PY{+w}{ }\PY{o}{=}\PY{+w}{ }\PY{l+s+s2}{\PYZdq{}}\PY{l+s+s2}{0.1.0}\PY{l+s+s2}{\PYZdq{}}
\PY{n}{channels}\PY{+w}{ }\PY{o}{=}\PY{+w}{ }\PY{p}{[}\PY{l+s+s2}{\PYZdq{}}\PY{l+s+s2}{robostack\PYZhy{}jazzy}\PY{l+s+s2}{\PYZdq{}}\PY{p}{,}\PY{+w}{ }\PY{l+s+s2}{\PYZdq{}}\PY{l+s+s2}{conda\PYZhy{}forge}\PY{l+s+s2}{\PYZdq{}}\PY{p}{]}
\PY{n}{platforms}\PY{+w}{ }\PY{o}{=}\PY{+w}{ }\PY{p}{[}\PY{l+s+s2}{\PYZdq{}}\PY{l+s+s2}{linux\PYZhy{}64}\PY{l+s+s2}{\PYZdq{}}\PY{p}{]}

\PY{k}{[}\PY{k}{dependencies}\PY{k}{]}
\PY{n}{ros\PYZhy{}jazzy\PYZhy{}desktop}\PY{+w}{ }\PY{o}{=}\PY{+w}{ }\PY{l+s+s2}{\PYZdq{}}\PY{l+s+s2}{*}\PY{l+s+s2}{\PYZdq{}}
\PY{n}{colcon\PYZhy{}common\PYZhy{}extensions}\PY{+w}{ }\PY{o}{=}\PY{+w}{ }\PY{l+s+s2}{\PYZdq{}}\PY{l+s+s2}{*}\PY{l+s+s2}{\PYZdq{}}

\PY{k}{[}\PY{k}{tasks}\PY{k}{]}
\PY{n}{build}\PY{+w}{ }\PY{o}{=}\PY{+w}{ }\PY{p}{\PYZob{}}\PY{+w}{ }\PY{n}{cmd}\PY{+w}{ }\PY{p}{=}\PY{+w}{ }\PY{l+s+s2}{\PYZdq{}}\PY{l+s+s2}{colcon build \PYZhy{}\PYZhy{}symlink\PYZhy{}install}\PY{l+s+s2}{\PYZdq{}}\PY{p}{,}\PY{+w}{ }\PY{n}{inputs}\PY{+w}{ }\PY{p}{=}\PY{+w}{ }\PY{p}{[}\PY{l+s+s2}{\PYZdq{}}\PY{l+s+s2}{src}\PY{l+s+s2}{\PYZdq{}}\PY{p}{]}\PY{+w}{ }\PY{p}{\PYZcb{}}
\PY{n}{sim}\PY{+w}{ }\PY{o}{=}\PY{+w}{ }\PY{l+s+s2}{\PYZdq{}}\PY{l+s+s2}{ros2 run turtlesim turtlesim\PYZus{}node}\PY{l+s+s2}{\PYZdq{}}

\PY{k}{[}\PY{k}{activation}\PY{k}{]}
\PY{n}{scripts}\PY{+w}{ }\PY{o}{=}\PY{+w}{ }\PY{p}{[}\PY{l+s+s2}{\PYZdq{}}\PY{l+s+s2}{install/setup.sh}\PY{l+s+s2}{\PYZdq{}}\PY{p}{]}
\end{Verbatim}
\end{filelisting}
\end{frame}

\begin{frame}{Secciones del manifiesto}
\begin{itemize}
    \item \lstinline|[workspace]| — metadatos + \lstinline|channels| + \lstinline|platforms|.
    \item \lstinline|[dependencies]| (conda) / \lstinline|[pypi-dependencies]| (PyPI, resuelto con \texttt{uv}) — se pueden mezclar.
    \item \lstinline|[tasks]| — comandos del proyecto.
    \item \lstinline|[activation]| — scripts/env vars que se cargan automáticamente al activar el entorno.
    \item \lstinline|[feature.*]| + \lstinline|[environments]| — varios entornos combinables en el mismo repo.
    \item \lstinline|[target.<platform>.*]| — overrides específicos de plataforma.
\end{itemize}
\vfill
\small También puede vivir embebido en \lstinline|pyproject.toml| bajo \lstinline|[tool.pixi.*]|.
\end{frame}

\begin{frame}[fragile]{Flujo de trabajo típico}
\begin{terminal}
pixi init                # crea el workspace (pixi.toml)
pixi add <paquete>       # añade dependencia, resuelve, instala
pixi install             # instala desde el lockfile existente
pixi run <comando>       # ejecuta dentro del entorno, sin "activate"
pixi shell               # abre un shell interactivo con el entorno activado
pixi task add / list     # gestiona tasks
pixi lock / pixi update  # gestiona el lockfile
pixi global install      # instala herramientas CLI globales, aisladas
\end{terminal}
\end{frame}

\begin{frame}{Lockfiles: el porqué de la reproducibilidad}
\begin{itemize}
    \item \lstinline|pixi.lock| fija versión exacta, build hash y \textbf{sha256} de cada paquete (conda y PyPI), \textbf{por cada plataforma declarada}.
    \item Se actualiza automáticamente con \lstinline|pixi add| / \lstinline|pixi install| / \lstinline|pixi update|.
    \item \lstinline|pixi install --locked| $\rightarrow$ falla si el lockfile no está sincronizado (ideal para CI).
    \item \lstinline|pixi install --frozen| $\rightarrow$ instala exactamente lo que dice el lockfile.
\end{itemize}
\vfill
\begin{block}{Mensaje central}
\lstinline|pixi.lock| se commitea a git: todo el equipo, CI y hasta el propio robot instalan exactamente los mismos binarios, byte a byte.
\end{block}
\end{frame}

\begin{frame}{Tasks y entornos multiplataforma (aperitivo)}
\begin{itemize}
    \item Las tasks pueden encadenarse (\lstinline|depends-on|), cachear por \lstinline|inputs|/\lstinline|outputs| (no recompila si nada cambió), y tener parámetros con Jinja.
    \item \lstinline|platforms = [...]| permite resolver dependencias distintas por SO/arquitectura en el mismo \lstinline|pixi.lock| — incluso requisitos ``ricos'' como CUDA o versión mínima de glibc.
\end{itemize}
\vfill
\centering \small Todo esto se practica en el siguiente bloque.
\end{frame}

\begin{frame}{Ventajas concretas para ROS~2}
\centering
\begin{tabular}{p{5.2cm}p{6.2cm}}
\toprule
\textbf{Problema} & \textbf{Cómo lo resuelve Pixi} \\
\midrule
ROS atado a versión de Ubuntu & Paquetes conda multiplataforma, sin depender de \texttt{apt} \\
Sin reproducibilidad real & \lstinline|pixi.lock| fija todo, byte a byte, por plataforma \\
Overlays manuales y frágiles & \lstinline|[activation]| automatiza el \texttt{source} en cada \lstinline|pixi shell|/\lstinline|run| \\
Varias distros ROS a la vez & \lstinline|[environments]| — Humble y Jazzy en paralelo, mismo repo \\
\bottomrule
\end{tabular}
\end{frame}

\begin{frame}{Un dato honesto}
\begin{alertblock}{RoboStack es hoy Tier 3}
En la clasificación oficial de plataformas de ROS: soporte \textbf{comunitario}, no de Open Robotics.
\end{alertblock}
\vfill
Pero con adopción creciente, y comunicación de que formará parte de la estrategia de accesibilidad de Open Robotics.
\end{frame}

%=================================================
\section[Práctica: Pixi básico]{Instalación y uso básico de Pixi}
%=================================================

\practicaframe{Instalación y uso básico de Pixi}{30 minutos}

\begin{frame}{Objetivo de este bloque}
\begin{itemize}
    \item Instalar Pixi y crear el primer workspace.
    \item Añadir dependencias, ejecutar comandos, definir tasks encadenadas con caché.
    \item Practicar múltiples entornos en el mismo proyecto.
\end{itemize}
\vfill
\begin{alertblock}{Todavía sin ROS}
Usamos ejemplos ligeros en Python para dominar la mecánica de Pixi sin mezclar curvas de aprendizaje. ROS llega en el bloque 5.
\end{alertblock}
\end{frame}

\begin{frame}[fragile]{Ejercicio 1 — Instalación y primer workspace}
\begin{terminal}
curl -fsSL https://pixi.sh/install.sh | sh
# reiniciar la terminal
pixi --version
\end{terminal}
\begin{terminal}
pixi init pixi-demo
cd pixi-demo
cat pixi.toml
ls -la
\end{terminal}
\vfill
\textbf{Checkpoint}: \lstinline|pixi.toml| recién generado con \lstinline|[workspace]|, además de \lstinline|.gitignore| y \lstinline|.gitattributes|.
\end{frame}

\begin{frame}[fragile]{Ejercicio 2 — Dependencias y ejecución}
\begin{terminal}
pixi add python=3.11 numpy rich
pixi run python -c "import numpy; print(numpy.__version__)"
pixi list
pixi tree numpy
\end{terminal}
\vfill
\textbf{Checkpoint}: \lstinline|pixi list| muestra versiones exactas; \lstinline|pixi tree| muestra dependencias transitivas. \lstinline|pixi.lock| ha cambiado.
\end{frame}

\preguntaframe{¿\lstinline|pixi add --pypi requests| y \lstinline|pixi add requests| instalan lo mismo?}

\begin{frame}[fragile]{Ejercicio 3 — Tasks encadenadas y con caché}
\begin{filelisting}
\begin{Verbatim}[commandchars=\\\{\}]
\PY{k}{[}\PY{k}{tasks}\PY{k}{]}
\PY{n}{hello}\PY{+w}{ }\PY{o}{=}\PY{+w}{ }\PY{l+s+s2}{\PYZdq{}}\PY{l+s+s2}{python \PYZhy{}c }\PY{l+s+se}{\PYZbs{}\PYZdq{}}\PY{l+s+s2}{print(\PYZsq{}hola workshop\PYZsq{})}\PY{l+s+se}{\PYZbs{}\PYZdq{}}\PY{l+s+s2}{\PYZdq{}}
\PY{n}{build}\PY{+w}{ }\PY{o}{=}\PY{+w}{ }\PY{p}{\PYZob{}}\PY{+w}{ }\PY{n}{cmd}\PY{+w}{ }\PY{p}{=}\PY{+w}{ }\PY{l+s+s2}{\PYZdq{}}\PY{l+s+s2}{python \PYZhy{}m py\PYZus{}compile main.py}\PY{l+s+s2}{\PYZdq{}}\PY{p}{,}
\PY{+w}{          }\PY{n}{depends\PYZhy{}on}\PY{+w}{ }\PY{o}{=}\PY{+w}{ }\PY{p}{[}\PY{l+s+s2}{\PYZdq{}}\PY{l+s+s2}{hello}\PY{l+s+s2}{\PYZdq{}}\PY{p}{]}\PY{err}{,}\PY{+w}{ }\PY{n}{inputs}\PY{+w}{ }\PY{o}{=}\PY{+w}{ }\PY{p}{[}\PY{l+s+s2}{\PYZdq{}}\PY{l+s+s2}{main.py}\PY{l+s+s2}{\PYZdq{}}\PY{p}{]}\PY{+w}{ }\PY{err}{\PYZcb{}}
\end{Verbatim}
\end{filelisting}
\begin{terminal}
echo "print('soy main.py')" > main.py
pixi run build
pixi task list
\end{terminal}
\vfill
\textbf{Reto}: ejecutar \lstinline|pixi run build| dos veces sin tocar \lstinline|main.py|. ¿Qué esperáis que cambie?
\end{frame}

\begin{frame}[fragile]{Ejercicio 4 — Multi-entorno con features}
\begin{filelisting}
\begin{Verbatim}[commandchars=\\\{\}]
\PY{k}{[}\PY{k}{dependencies}\PY{k}{]}\PY{+w}{  }\PY{c+c1}{\PYZsh{} feature \PYZdq{}default\PYZdq{}: ¡ya trae python 3.11!}
\PY{n}{python}\PY{+w}{ }\PY{o}{=}\PY{+w}{ }\PY{l+s+s2}{\PYZdq{}}\PY{l+s+s2}{3.11.*}\PY{l+s+s2}{\PYZdq{}}

\PY{k}{[}\PY{k}{feature}\PY{k}{.}\PY{k}{py310}\PY{k}{.}\PY{k}{dependencies}\PY{k}{]}
\PY{n}{python}\PY{+w}{ }\PY{o}{=}\PY{+w}{ }\PY{l+s+s2}{\PYZdq{}}\PY{l+s+s2}{3.10.*}\PY{l+s+s2}{\PYZdq{}}
\PY{k}{[}\PY{k}{feature}\PY{k}{.}\PY{k}{py311}\PY{k}{.}\PY{k}{dependencies}\PY{k}{]}
\PY{n}{python}\PY{+w}{ }\PY{o}{=}\PY{+w}{ }\PY{l+s+s2}{\PYZdq{}}\PY{l+s+s2}{3.11.*}\PY{l+s+s2}{\PYZdq{}}

\PY{k}{[}\PY{k}{environments}\PY{k}{]}
\PY{n}{py310}\PY{+w}{ }\PY{o}{=}\PY{+w}{ }\PY{p}{\PYZob{}}\PY{+w}{ }\PY{n}{features}\PY{+w}{ }\PY{p}{=}\PY{+w}{ }\PY{p}{[}\PY{l+s+s2}{\PYZdq{}}\PY{l+s+s2}{py310}\PY{l+s+s2}{\PYZdq{}}\PY{p}{]}\PY{p}{,}\PY{+w}{ }\PY{n}{no\PYZhy{}default\PYZhy{}feature}\PY{+w}{ }\PY{p}{=}\PY{+w}{ }\PY{k+kc}{true}\PY{+w}{ }\PY{p}{\PYZcb{}}
\PY{n}{py311}\PY{+w}{ }\PY{o}{=}\PY{+w}{ }\PY{p}{\PYZob{}}\PY{+w}{ }\PY{n}{features}\PY{+w}{ }\PY{p}{=}\PY{+w}{ }\PY{p}{[}\PY{l+s+s2}{\PYZdq{}}\PY{l+s+s2}{py311}\PY{l+s+s2}{\PYZdq{}}\PY{p}{]}\PY{p}{,}\PY{+w}{ }\PY{n}{no\PYZhy{}default\PYZhy{}feature}\PY{+w}{ }\PY{p}{=}\PY{+w}{ }\PY{k+kc}{true}\PY{+w}{ }\PY{p}{\PYZcb{}}
\end{Verbatim}
\end{filelisting}
\begin{terminal}
pixi install
pixi run -e py310 python --version
pixi run -e py311 python --version
\end{terminal}
\small \textbf{Checkpoint}: dos versiones de Python distintas, mismo repo, sin virtualenvs.\\
\small \textbf{Trampa}: cada entorno de \lstinline|[environments]| incluye la feature \lstinline|default| a menos que uses \lstinline|no-default-feature = true| — si no, \lstinline|python 3.11| (default) choca con \lstinline|python 3.10| (py310) y falla el solve.
\end{frame}

\begin{frame}{Cierre del bloque}
\centering
\Large \lstinline|init| \quad \lstinline|add| \quad \lstinline|run| \quad \lstinline|shell| \quad \lstinline|task| \quad \lstinline|environments|
\vfill
\normalsize
Esto —dos versiones de Python conviviendo en el mismo proyecto— es exactamente el mecanismo que usaremos para tener ROS~2 Humble y Jazzy en el mismo repo sin pisarse.
\end{frame}

%=================================================
\section[¿Qué es RoboStack?]{¿Qué es RoboStack?}
%=================================================

\begin{frame}{¿Qué es RoboStack?}
\begin{itemize}
    \item Empaqueta distribuciones completas de ROS~1/2 como \textbf{paquetes conda}.
    \item Convención de nombres: \lstinline|ros-<distro>-<paquete>|, p.\,ej. \lstinline|ros-jazzy-desktop|, \lstinline|ros-jazzy-rclcpp|.
    \item Permite \lstinline|pixi add ros-jazzy-desktop| igual que se añadiría \texttt{numpy}.
\end{itemize}
\end{frame}

\begin{frame}{Origen y motivación}
\begin{itemize}
    \item Paper académico: \emph{``A RoboStack Tutorial''}, Fischer et al., IEEE Robotics \& Automation Magazine, 2021 (arXiv:2104.12910).
    \item Wolf Vollprecht (creador de mamba, cofundador de prefix.dev) y Tobias Fischer, entre otros.
    \item Motivación original: mezclar ROS con el ecosistema de ciencia de datos (PyTorch, Jupyter, OpenCV) sin pelearse con versiones fijas de \texttt{apt}, sin \texttt{sudo}, y en macOS/Windows además de Linux.
\end{itemize}
\end{frame}

\begin{frame}{Canales conda}
\centering
\begin{tabular}{p{4cm}p{7cm}}
\toprule
\textbf{Canal} & \textbf{Para qué sirve} \\
\midrule
\lstinline|robostack-<distro>| & \textbf{Recomendado para uso normal.} Uno por distro ROS. \\
\lstinline|robostack-staging| & Interno/histórico, proceso de build de los mantenedores. \\
\lstinline|conda-forge| & Dependencias no-ROS (compiladores, Boost, OpenCV, Python\ldots). \\
\bottomrule
\end{tabular}
\vfill
\begin{alertblock}{Regla simple}
Para uso normal, siempre \lstinline|robostack-<vuestra-distro>| + \lstinline|conda-forge|. Nada más.
\end{alertblock}
\end{frame}

\begin{frame}{Compatibilidad multiplataforma}
\begin{itemize}
    \item Linux: \lstinline|linux-64|, \lstinline|linux-aarch64|
    \item macOS: \lstinline|osx-64|, \lstinline|osx-arm64| (sí, ROS~2 en Apple Silicon)
    \item Windows: \lstinline|win-64|
\end{itemize}
\vfill
CI propio por plataforma en GitHub Actions.
\end{frame}

\begin{frame}{Aviso importante}
\begin{alertblock}{\texttt{rosdep} no funciona dentro de Pixi/RoboStack}
Se salta la resolución de Pixi llamando directamente a \texttt{apt}/\texttt{pip} del sistema.
\end{alertblock}
\vfill
\begin{itemize}
    \item Las dependencias de un paquete propio hay que traducirlas a mano: \lstinline|pixi add ros-<distro>-<paquete>|.
    \item O usar la herramienta comunitaria \textbf{pixi-ros}, que lee los \lstinline|package.xml| del workspace y genera el \lstinline|pixi.toml|.
\end{itemize}
\end{frame}

\begin{frame}{Relación con Pixi}
\begin{center}
\Large RoboStack no sustituye a Pixi.\\
Es ``solo'' un conjunto de canales conda más.
\end{center}
\vfill
Todo lo del bloque anterior (\lstinline|add|, \lstinline|run|, \lstinline|shell|, tasks, \lstinline|[environments]|, \lstinline|[activation]|) se aplica sin cambios.
\end{frame}

\preguntaframe{Con lo que sabéis de Pixi, ¿cómo creéis que se instalaría ROS~2 Jazzy Desktop?}

%=================================================
\section[Práctica: RoboStack]{Uso básico de RoboStack para compilar un workspace ROS~2}
%=================================================

\practicaframe{RoboStack + colcon}{20 minutos}

\begin{frame}[fragile]{Ejercicio 1 — Entorno ROS~2 en dos comandos}
\begin{terminal}
pixi init taller_ros2 -c robostack-jazzy -c conda-forge
cd taller_ros2
pixi add ros-jazzy-desktop
pixi run ros2 topic list
\end{terminal}
\begin{terminal}
# Terminal 1
pixi run ros2 run demo_nodes_cpp talker
# Terminal 2
pixi run ros2 run demo_nodes_cpp listener
\end{terminal}
\vfill
\textbf{Checkpoint}: el listener recibe los mensajes del talker. Sin \texttt{sudo}, sin tocar el sistema.
\end{frame}

\begin{frame}[fragile]{Ejercicio 2 — Paquete propio + colcon}
\begin{terminal}
pixi add colcon-common-extensions "setuptools<=58.2.0" "empy==3.3.4"
mkdir src
pixi run ros2 pkg create --build-type ament_python \
  --destination-directory src --node-name talker_custom mi_paquete
pixi run colcon build --symlink-install
\end{terminal}
\vfill
\textbf{Checkpoint}: aparecen \lstinline|build/|, \lstinline|install/|, \lstinline|log/| junto a \lstinline|src/|.
\end{frame}

\begin{frame}[fragile]{El overlay, automatizado}
\begin{filelisting}
\begin{Verbatim}[commandchars=\\\{\}]
\PY{c+c1}{\PYZsh{} añadir a pixi.toml}
\PY{k}{[}\PY{k}{activation}\PY{k}{]}
\PY{n}{scripts}\PY{+w}{ }\PY{o}{=}\PY{+w}{ }\PY{p}{[}\PY{l+s+s2}{\PYZdq{}}\PY{l+s+s2}{install/setup.sh}\PY{l+s+s2}{\PYZdq{}}\PY{p}{]}
\end{Verbatim}
\end{filelisting}
\begin{terminal}
pixi shell
echo $AMENT_PREFIX_PATH   # install/ del workspace, primero
ros2 run mi_paquete talker_custom
exit
\end{terminal}
\vfill
\textbf{Checkpoint}: \lstinline|$AMENT_PREFIX_PATH| lista \lstinline|.../taller_ros2/install| como primera entrada. Overlay funcionando, sin \lstinline|source install/setup.bash| a mano.
\end{frame}

\preguntaframe{¿Qué pasaría si borro \texttt{install/} pero dejo el \texttt{[activation]} en el \texttt{pixi.toml}?}

\begin{frame}[fragile]{Ejercicio 3 — Tasks con caché encadenadas}
\begin{filelisting}
\begin{Verbatim}[commandchars=\\\{\}]
\PY{k}{[}\PY{k}{tasks}\PY{k}{]}
\PY{n}{build}\PY{+w}{  }\PY{o}{=}\PY{+w}{ }\PY{p}{\PYZob{}}\PY{+w}{ }\PY{n}{cmd}\PY{+w}{ }\PY{p}{=}\PY{+w}{ }\PY{l+s+s2}{\PYZdq{}}\PY{l+s+s2}{colcon build \PYZhy{}\PYZhy{}symlink\PYZhy{}install}\PY{l+s+s2}{\PYZdq{}}\PY{p}{,}\PY{+w}{ }\PY{n}{inputs}\PY{+w}{ }\PY{p}{=}\PY{+w}{ }\PY{p}{[}\PY{l+s+s2}{\PYZdq{}}\PY{l+s+s2}{src}\PY{l+s+s2}{\PYZdq{}}\PY{p}{]}\PY{+w}{ }\PY{p}{\PYZcb{}}
\PY{n}{run\PYZus{}talker}\PY{+w}{ }\PY{o}{=}\PY{+w}{ }\PY{p}{\PYZob{}}\PY{+w}{ }\PY{n}{cmd}\PY{+w}{ }\PY{p}{=}\PY{+w}{ }\PY{l+s+s2}{\PYZdq{}}\PY{l+s+s2}{ros2 run mi\PYZus{}paquete talker\PYZus{}custom}\PY{l+s+s2}{\PYZdq{}}\PY{p}{,}
\PY{+w}{               }\PY{n}{depends\PYZhy{}on}\PY{+w}{ }\PY{o}{=}\PY{+w}{ }\PY{p}{[}\PY{l+s+s2}{\PYZdq{}}\PY{l+s+s2}{build}\PY{l+s+s2}{\PYZdq{}}\PY{p}{]}\PY{+w}{ }\PY{err}{\PYZcb{}}
\end{Verbatim}
\end{filelisting}
\begin{terminal}
pixi run run_talker
pixi run run_talker   # 2a vez: ¿recompila?
# editar talker_custom.py y repetir
pixi run run_talker   # 3a vez: ¿y ahora?
\end{terminal}
\end{frame}

\begin{frame}{Ejercicio 4 — Un workspace real: Nav2 con \texttt{pixi-ros}}
\begin{itemize}
    \item Hasta ahora, un paquete de juguete. ¿Y un workspace real y grande?
    \item Y de paso, resolvemos lo que dejamos abierto en el bloque anterior: si \texttt{rosdep} no funciona, ¿cómo instalo las dependencias de un \lstinline|package.xml| real sin hacerlo a mano?
\end{itemize}
\vfill
\begin{block}{\texttt{pixi-ros}}
Herramienta de la comunidad Pixi: escanea los \lstinline|package.xml| del workspace y rellena el \lstinline|pixi.toml| con las dependencias ya resueltas a paquetes conda de RoboStack. El \texttt{rosdep install} moderno, pero reproducible.
\end{block}
\vfill
\begin{alertblock}{Aviso de timing}
Nav2 son más de 80 paquetes C++. El instructor lanza el build al empezar el bloque anterior — aquí mostramos el resultado en vivo.
\end{alertblock}
\end{frame}

\begin{frame}[fragile]{Preparar el workspace}
\begin{terminal}
pixi global install pixi-ros
\end{terminal}
\begin{terminal}
mkdir -p nav2_ws/src && cd nav2_ws
git clone -b jazzy https://github.com/ros-navigation/navigation2.git \
  src/navigation2
\end{terminal}
\end{frame}

\begin{frame}[fragile]{Rellenar dependencias automáticamente}
\begin{terminal}
pixi-ros init --distro jazzy --platform linux-64
\end{terminal}
\begin{itemize}
    \item Escanea todos los \lstinline|package.xml| (\lstinline|nav2_bt_navigator|, \lstinline|nav2_costmap_2d|, \lstinline|nav2_planner|\ldots).
    \item Resuelve cada dependencia por prioridad: workspace $\rightarrow$ ficheros de mapeo $\rightarrow$ índice de la distro ROS $\rightarrow$ conda-forge $\rightarrow$ canales extra.
    \item Muestra una tabla de validación con el origen de cada dependencia.
\end{itemize}
\vfill
\textbf{Checkpoint}: el \lstinline|pixi.toml| generado tiene decenas de \lstinline|ros-jazzy-*| + herramientas de build, y un bloque \lstinline|[tasks]| con \lstinline|build|/\lstinline|test| ya creados (\texttt{colcon} por debajo). Es idempotente.
\end{frame}

\begin{frame}[fragile]{Ajustar dependencias}
\begin{terminal}
pixi add gcc_linux-64=13.3.0 gxx_linux-64=13.3.0
pixi add libboost=1.88.0 libboost-devel=1.88.0
pixi add ros-jazzy-rviz2
\end{terminal}
\begin{itemize}
    \item Por defecto, pixi-ros usa una versión de GCC y de Boost demasiado moderna para ROS 2 Jazzy. Para compilar nav2 necesitamos fijarlas a las versiones usadas en Ubuntu 24.04.
    \item Rviz no se añade automáticamente por pixi-ros porque no es una dependencia de build, pero se usará en ejecución.
\end{itemize}
\end{frame}

\begin{frame}[fragile]{Ajustar dependencias}
\begin{alertblock}{Error de sintaxis en TOML generado por \texttt{pixi-ros init}}
En algunas versiones, la dependencia fijada de \texttt{xtensor} se escribe como una clave TOML inválida.
\end{alertblock}
\begin{filelisting}
\begin{Verbatim}[commandchars=\\\{\}]
\PY{k}{[}\PY{k}{dependencies}\PY{k}{]}
\PY{+w}{...}
\PY{l+s+s2}{\PYZdq{}}\PY{l+s+s2}{xtensor ==0.24.7}\PY{l+s+s2}{\PYZdq{}}\PY{+w}{ }\PY{o}{=}\PY{+w}{ }\PY{l+s+s2}{\PYZdq{}}\PY{l+s+s2}{*}\PY{l+s+s2}{\PYZdq{}}
\end{Verbatim}
\end{filelisting}
\begin{center}
\Large $\Downarrow$ corregir en \texttt{pixi.toml}
\end{center}
\begin{filelisting}
\begin{Verbatim}[commandchars=\\\{\}]
\PY{k}{[}\PY{k}{dependencies}\PY{k}{]}
\PY{+w}{...}
\PY{n}{xtensor}\PY{+w}{ }\PY{o}{=}\PY{+w}{ }\PY{l+s+s2}{\PYZdq{}}\PY{l+s+s2}{==0.24.7}\PY{l+s+s2}{\PYZdq{}}
\end{Verbatim}
\end{filelisting}
\end{frame}

\begin{frame}[fragile]{Compilar y verificar}
\begin{terminal}
pixi install
pixi run build
\end{terminal}
\begin{terminal}
pixi shell
ros2 pkg list | grep nav2
ros2 pkg executables nav2_costmap_2d
\end{terminal}
\vfill
\textbf{Checkpoint}: decenas de paquetes \lstinline|nav2_*|, compilados dentro de este entorno Pixi, sin tocar el sistema ni usar \texttt{rosdep}.
\end{frame}

\begin{frame}[fragile]{El cierre real: simulación + RViz2}
\lstinline|nav2_bringup| declara en su \lstinline|package.xml|: \lstinline|ros_gz_bridge|, \lstinline|ros_gz_sim|, \lstinline|xacro|, \lstinline|slam_toolbox|, \lstinline|nav2_minimal_tb3_sim|.
\vfill
\begin{block}{}
Todo ya resuelto por \texttt{pixi-ros} como binarios conda de RoboStack. No hace falta clonar ningún repo más.
\end{block}
\begin{terminal}
ros2 launch nav2_bringup tb3_simulation_launch.py
\end{terminal}
\vfill
Levanta \textbf{Gazebo Sim} (\lstinline|gz sim|) con el robot spawneado, toda la pila de Nav2, y \textbf{RViz2} con la vista por defecto.
\end{frame}

\begin{frame}{Mandar al robot a navegar, en RViz2}
\begin{enumerate}
    \item Botón \textbf{2D Pose Estimate}: click y arrastra sobre el mapa en la posición/orientación real del robot (localizarlo para AMCL).
    \item Botón \textbf{Nav2 Goal}: click y arrastra en el punto de destino, con la orientación final deseada.
    \item El robot planifica una ruta y se mueve solo hasta el objetivo, esquivando obstáculos.
\end{enumerate}
\vfill
\textbf{Checkpoint}: el robot navega autónomamente en la simulación — Gazebo, RViz2 y Nav2 completo, dentro del mismo entorno Pixi, sin \texttt{sudo}, sin \texttt{apt}, sin Docker.
\vfill
\begin{alertblock}{Aviso honesto: GPU/OpenGL}
\lstinline|gz sim| necesita una pila gráfica funcional. En VMs/WSL2 sin passthrough de GPU puede fallar o ir lento — probadlo antes del taller. Fallback: \lstinline|LIBGL_ALWAYS_SOFTWARE=1| delante del comando, o hacerlo como demo proyectada del instructor.
\end{alertblock}
\end{frame}

\begin{frame}[fragile]{Ejercicio 5.5 (opcional) — Varias distros ROS~2 a la vez (1/2)}
Mismo patrón que el ejercicio 3.4 (Python), pero cada distro RoboStack vive en su propio canal: no se pueden mezclar \lstinline|robostack-jazzy| y \lstinline|robostack-humble| en la misma resolución.
\begin{filelisting}
\begin{Verbatim}[commandchars=\\\{\}]
\PY{k}{[}\PY{k}{feature}\PY{k}{.}\PY{k}{jazzy}\PY{k}{]}
\PY{n}{channels}\PY{+w}{ }\PY{o}{=}\PY{+w}{ }\PY{p}{[}\PY{l+s+s2}{\PYZdq{}}\PY{l+s+s2}{robostack\PYZhy{}jazzy}\PY{l+s+s2}{\PYZdq{}}\PY{p}{,}\PY{+w}{ }\PY{l+s+s2}{\PYZdq{}}\PY{l+s+s2}{conda\PYZhy{}forge}\PY{l+s+s2}{\PYZdq{}}\PY{p}{]}
\PY{k}{[}\PY{k}{feature}\PY{k}{.}\PY{k}{jazzy}\PY{k}{.}\PY{k}{dependencies}\PY{k}{]}
\PY{n}{ros\PYZhy{}jazzy\PYZhy{}desktop}\PY{+w}{ }\PY{o}{=}\PY{+w}{ }\PY{l+s+s2}{\PYZdq{}}\PY{l+s+s2}{*}\PY{l+s+s2}{\PYZdq{}}

\PY{k}{[}\PY{k}{feature}\PY{k}{.}\PY{k}{humble}\PY{k}{]}
\PY{n}{channels}\PY{+w}{ }\PY{o}{=}\PY{+w}{ }\PY{p}{[}\PY{l+s+s2}{\PYZdq{}}\PY{l+s+s2}{robostack\PYZhy{}humble}\PY{l+s+s2}{\PYZdq{}}\PY{p}{,}\PY{+w}{ }\PY{l+s+s2}{\PYZdq{}}\PY{l+s+s2}{conda\PYZhy{}forge}\PY{l+s+s2}{\PYZdq{}}\PY{p}{]}
\PY{k}{[}\PY{k}{feature}\PY{k}{.}\PY{k}{humble}\PY{k}{.}\PY{k}{dependencies}\PY{k}{]}
\PY{n}{ros\PYZhy{}humble\PYZhy{}desktop}\PY{+w}{ }\PY{o}{=}\PY{+w}{ }\PY{l+s+s2}{\PYZdq{}}\PY{l+s+s2}{*}\PY{l+s+s2}{\PYZdq{}}

\PY{k}{[}\PY{k}{environments}\PY{k}{]}
\PY{n}{jazzy}\PY{+w}{  }\PY{o}{=}\PY{+w}{ }\PY{p}{[}\PY{l+s+s2}{\PYZdq{}}\PY{l+s+s2}{jazzy}\PY{l+s+s2}{\PYZdq{}}\PY{p}{]}
\PY{n}{humble}\PY{+w}{ }\PY{o}{=}\PY{+w}{ }\PY{p}{[}\PY{l+s+s2}{\PYZdq{}}\PY{l+s+s2}{humble}\PY{l+s+s2}{\PYZdq{}}\PY{p}{]}
\end{Verbatim}
\end{filelisting}
\end{frame}

\begin{frame}[fragile]{Ejercicio 5.5 (opcional) — Varias distros ROS~2 a la vez (2/2)}
\begin{terminal}
pixi install
pixi run -e jazzy  ros2 --version
pixi run -e humble ros2 --version
\end{terminal}
\vfill
\textbf{Checkpoint}: cada comando reporta la versión de \lstinline|ros2| de su propia distro — dos ROS~2 completos, en el mismo repositorio, sin pisarse.
\end{frame}

\begin{frame}{Dos ROS a la vez, sin pisarse}
\begin{center}
\Large Imposible o muy incómodo con \texttt{apt}.\\[0.3em]
Con Pixi, un \lstinline|[environments]| más.
\end{center}
\vfill
\normalsize Exactamente el mismo mecanismo del bloque 3 con Python — solo cambia qué metemos dentro de cada \emph{feature}.
\end{frame}

%=================================================
\section[Infraestructura RoboStack]{Infraestructura de RoboStack}
%=================================================

\begin{frame}{De \texttt{package.xml} a paquete conda}
En el bloque anterior usamos paquetes ya construidos. ¿Cómo se construyen?
\begin{itemize}
    \item \textbf{vinca}: lee el \emph{rosdistro} oficial (metadata + \lstinline|package.xml| de cada paquete) y un fichero \lstinline|vinca.yaml|, y genera recetas en formato \textbf{rattler-build} (\lstinline|recipe.yaml|, esquema \texttt{prefix-dev/recipe-format}).
    \item \textbf{rattler-build}: el motor que compila cada receta y produce el \lstinline|.conda| final. Sucesor moderno de \texttt{conda-build}/\texttt{boa}.
\end{itemize}
\end{frame}

\begin{frame}{Recetas y build farms}
\begin{itemize}
    \item Cada repo de distro (\texttt{RoboStack/ros-jazzy}, \texttt{ros-humble}, \texttt{ros-kilted}\ldots) contiene la \emph{configuración} (\lstinline|vinca.yaml|, \lstinline|robostack.yaml|, \lstinline|patches/|, \lstinline|conda_build_config.yaml|) — pero \textbf{no} recetas estáticas: se regeneran en cada build.
    \item CI en \textbf{GitHub Actions}, matriz por plataforma (\lstinline|linux-64|, \lstinline|linux-aarch64|, \lstinline|osx-64|, \lstinline|osx-arm64|, \lstinline|win-64|), con caché de builds y generación dinámica de workflows.
\end{itemize}
\end{frame}

\begin{frame}{Relación con conda-forge y contribución}
\begin{itemize}
    \item RoboStack se apoya en \textbf{conda-forge} para todas las dependencias no-ROS (compiladores, Boost, PCL, OpenCV, Python\ldots).
    \item Contribuir un paquete nuevo, en general: PR de una línea en \lstinline|vinca.yaml| $\rightarrow$ CI intenta construirlo en todas las plataformas $\rightarrow$ si falla, parche local guardado en \lstinline|patches/| $\rightarrow$ rebuild para validar.
\end{itemize}
\end{frame}

\begin{frame}
\begin{center}
\resizebox{0.92\linewidth}{!}{%
\begin{tikzpicture}[node distance=1.1cm, every node/.style={font=\small}]
\node[draw, rounded corners, fill=blue!10, minimum width=3.2cm, minimum height=0.9cm] (pkgxml) {\lstinline|package.xml| (ROS)};
\node[draw, rounded corners, fill=orange!15, minimum width=3.2cm, minimum height=0.9cm, right=1.4cm of pkgxml] (vinca) {vinca};
\node[draw, rounded corners, fill=orange!15, minimum width=3.6cm, minimum height=0.9cm, right=1.4cm of vinca] (recipe) {\lstinline|recipe.yaml|};
\node[draw, rounded corners, fill=green!15, minimum width=3.4cm, minimum height=0.9cm, below=1.2cm of vinca] (rattler) {rattler-build};
\node[draw, rounded corners, fill=green!25, minimum width=3.4cm, minimum height=0.9cm, right=1.4cm of rattler] (conda) {paquete \lstinline|.conda|};
\draw[-Latex, thick] (pkgxml) -- (vinca);
\draw[-Latex, thick] (vinca) -- (recipe);
\draw[-Latex, thick] (recipe) |- (rattler);
\draw[-Latex, thick] (rattler) -- (conda);
\end{tikzpicture}%
}
\end{center}
\vfill
\centering \small Este es exactamente el pipeline que vais a recorrer, en pequeño, en el próximo bloque práctico.
\end{frame}

%=================================================
\section[Práctica: tu buildfarm]{Tu propia Buildfarm basada en RoboStack}
%=================================================

\practicaframe{Tu propia Buildfarm}{35 minutos}

\begin{frame}{Objetivo y nota de diseño}
\begin{itemize}
    \item Construir un paquete conda propio con \texttt{rattler-build}.
    \item Publicarlo en un canal de prefix.dev.
    \item Instalarlo desde un \lstinline|pixi.toml| propio.
    \item Integrarlo en un workspace ROS~2.
\end{itemize}
\vfill
\begin{alertblock}{No usamos Vinca hoy}
Generar recetas reales de ROS~2 completo tarda demasiado para 1h. Escribimos una receta \texttt{rattler-build} sencilla, con el mismo formato exacto que usa la buildfarm real.
\end{alertblock}
\end{frame}

\begin{frame}[fragile]{Fase A — Preparar el proyecto (1/3)}
\begin{terminal}
mkdir -p ~/taller-pixi/mi-paquete && cd ~/taller-pixi/mi-paquete
pixi init .
pixi add rattler-build rattler-index
mkdir -p src recipe
\end{terminal}
\begin{filelisting}
\begin{Verbatim}[commandchars=\\\{\}]
\PY{c+c1}{\PYZsh{} src/pyproject.toml}
\PY{k}{[}\PY{k}{project}\PY{k}{]}
\PY{n}{name}\PY{+w}{ }\PY{o}{=}\PY{+w}{ }\PY{l+s+s2}{\PYZdq{}}\PY{l+s+s2}{hola\PYZhy{}pichi}\PY{l+s+s2}{\PYZdq{}}
\PY{n}{version}\PY{+w}{ }\PY{o}{=}\PY{+w}{ }\PY{l+s+s2}{\PYZdq{}}\PY{l+s+s2}{1.0.0}\PY{l+s+s2}{\PYZdq{}}
\PY{k}{[}\PY{k}{project}\PY{k}{.}\PY{k}{scripts}\PY{k}{]}
\PY{n}{hola\PYZhy{}pichi}\PY{+w}{ }\PY{o}{=}\PY{+w}{ }\PY{l+s+s2}{\PYZdq{}}\PY{l+s+s2}{hola\PYZus{}pichi:main}\PY{l+s+s2}{\PYZdq{}}
\PY{k}{[}\PY{k}{build\PYZhy{}system}\PY{k}{]}
\PY{n}{requires}\PY{+w}{ }\PY{o}{=}\PY{+w}{ }\PY{p}{[}\PY{l+s+s2}{\PYZdq{}}\PY{l+s+s2}{setuptools}\PY{l+s+s2}{\PYZdq{}}\PY{p}{]}
\PY{n}{build\PYZhy{}backend}\PY{+w}{ }\PY{o}{=}\PY{+w}{ }\PY{l+s+s2}{\PYZdq{}}\PY{l+s+s2}{setuptools.build\PYZus{}meta}\PY{l+s+s2}{\PYZdq{}}
\end{Verbatim}
\end{filelisting}
\end{frame}

\begin{frame}[fragile]{Fase A — La receta \texttt{rattler-build} (2/3)}
\begin{filelisting}
\begin{Verbatim}[commandchars=\\\{\}]
\PY{c+c1}{\PYZsh{} recipe/recipe.yaml}
\PY{n+nt}{package}\PY{p}{:}
\PY{+w}{  }\PY{n+nt}{name}\PY{p}{:}\PY{+w}{ }\PY{l+lScalar+lScalarPlain}{hola\PYZhy{}pichi}
\PY{+w}{  }\PY{n+nt}{version}\PY{p}{:}\PY{+w}{ }\PY{l+s}{\PYZdq{}}\PY{l+s}{1.0.0}\PY{l+s}{\PYZdq{}}
\PY{n+nt}{source}\PY{p}{:}
\PY{+w}{  }\PY{n+nt}{path}\PY{p}{:}\PY{+w}{ }\PY{l+lScalar+lScalarPlain}{../src}
\PY{n+nt}{build}\PY{p}{:}
\PY{+w}{  }\PY{n+nt}{number}\PY{p}{:}\PY{+w}{ }\PY{l+lScalar+lScalarPlain}{0}
\PY{+w}{  }\PY{n+nt}{noarch}\PY{p}{:}\PY{+w}{ }\PY{l+lScalar+lScalarPlain}{python}
\PY{+w}{  }\PY{n+nt}{script}\PY{p}{:}\PY{+w}{ }\PY{l+lScalar+lScalarPlain}{python}\PY{l+lScalar+lScalarPlain}{ }\PY{l+lScalar+lScalarPlain}{\PYZhy{}m}\PY{l+lScalar+lScalarPlain}{ }\PY{l+lScalar+lScalarPlain}{pip}\PY{l+lScalar+lScalarPlain}{ }\PY{l+lScalar+lScalarPlain}{install}\PY{l+lScalar+lScalarPlain}{ }\PY{l+lScalar+lScalarPlain}{.}\PY{l+lScalar+lScalarPlain}{ }\PY{l+lScalar+lScalarPlain}{\PYZhy{}\PYZhy{}no\PYZhy{}deps}\PY{l+lScalar+lScalarPlain}{ }\PY{l+lScalar+lScalarPlain}{\PYZhy{}\PYZhy{}no\PYZhy{}build\PYZhy{}isolation}\PY{l+lScalar+lScalarPlain}{ }\PY{l+lScalar+lScalarPlain}{\PYZhy{}vv}
\PY{n+nt}{requirements}\PY{p}{:}
\PY{+w}{  }\PY{n+nt}{host}\PY{p}{:}\PY{+w}{ }\PY{p+pIndicator}{[}\PY{n+nv}{python}\PY{p+pIndicator}{,}\PY{+w}{ }\PY{n+nv}{pip}\PY{p+pIndicator}{]}
\PY{+w}{  }\PY{n+nt}{run}\PY{p}{:}\PY{+w}{ }\PY{p+pIndicator}{[}\PY{n+nv}{python}\PY{p+pIndicator}{]}
\PY{n+nt}{tests}\PY{p}{:}
\PY{+w}{  }\PY{p+pIndicator}{\PYZhy{}}\PY{+w}{ }\PY{n+nt}{script}\PY{p}{:}\PY{+w}{ }\PY{p+pIndicator}{[}\PY{n+nv}{hola\PYZhy{}pichi}\PY{p+pIndicator}{]}
\end{Verbatim}
\end{filelisting}
\end{frame}

\begin{frame}[fragile]{Fase A — Construir el paquete (3/3)}
\begin{terminal}
pixi run rattler-build build --recipe ./recipe/recipe.yaml -c conda-forge
ls output/noarch/*.conda
\end{terminal}
\vfill
\textbf{Checkpoint}: debe existir un fichero \lstinline|output/noarch/hola-pichi-1.0.0-*.conda|.
\vfill
\small Esta receta usa exactamente el mismo formato (\texttt{recipe-format} de prefix.dev) que genera Vinca automáticamente para cada paquete ROS~2.
\end{frame}

\begin{frame}[fragile]{Fase B — Publicar en prefix.dev}
\begin{terminal}
export PREFIX_API_KEY="pfx_..."

pixi run rattler-build upload prefix \
  --channel workshop-pichi-<nombre> \
  --skip-existing \
  output/noarch/*.conda
\end{terminal}
\vfill
\textbf{Checkpoint}: el paquete \texttt{hola-pichi} aparece listado en el canal, en prefix.dev.
\vfill
\begin{alertblock}{Canal público vs. privado}
Plan free: 1 canal privado disponible. Usamos uno público hoy por simplicidad — la gestión de permisos de canales privados no se cubre en este taller.
\end{alertblock}
\end{frame}

\begin{frame}[fragile]{Fase C — Instalar desde un \texttt{pixi.toml} propio}
\begin{filelisting}
\begin{Verbatim}[commandchars=\\\{\}]
\PY{c+c1}{\PYZsh{} pixi.toml}
\PY{k}{[}\PY{k}{workspace}\PY{k}{]}
\PY{n}{name}\PY{+w}{ }\PY{o}{=}\PY{+w}{ }\PY{l+s+s2}{\PYZdq{}}\PY{l+s+s2}{consumidor\PYZus{}ws}\PY{l+s+s2}{\PYZdq{}}
\PY{n}{channels}\PY{+w}{ }\PY{o}{=}\PY{+w}{ }\PY{p}{[}\PY{l+s+s2}{\PYZdq{}}\PY{l+s+s2}{https://prefix.dev/workshop\PYZhy{}pichi\PYZhy{}\PYZlt{}nombre\PYZgt{}}\PY{l+s+s2}{\PYZdq{}}\PY{p}{,}\PY{+w}{ }\PY{l+s+s2}{\PYZdq{}}\PY{l+s+s2}{conda\PYZhy{}forge}\PY{l+s+s2}{\PYZdq{}}\PY{p}{]}
\PY{n}{platforms}\PY{+w}{ }\PY{o}{=}\PY{+w}{ }\PY{p}{[}\PY{l+s+s2}{\PYZdq{}}\PY{l+s+s2}{linux\PYZhy{}64}\PY{l+s+s2}{\PYZdq{}}\PY{p}{]}

\PY{k}{[}\PY{k}{dependencies}\PY{k}{]}
\PY{n}{hola\PYZhy{}pichi}\PY{+w}{ }\PY{o}{=}\PY{+w}{ }\PY{l+s+s2}{\PYZdq{}}\PY{l+s+s2}{*}\PY{l+s+s2}{\PYZdq{}}
\end{Verbatim}
\end{filelisting}
\begin{terminal}
pixi install
pixi run hola-pichi
\end{terminal}
\vfill
\small \textbf{Troubleshooting}: si no aparece (\texttt{No candidates found}) $\rightarrow$ \lstinline|pixi clean cache --repodata -y|
\end{frame}

\begin{frame}[fragile]{Fase D — Integrarlo en un workspace ROS~2}
\begin{filelisting}
\begin{Verbatim}[commandchars=\\\{\}]
\PY{k}{[}\PY{k}{workspace}\PY{k}{]}
\PY{n}{channels}\PY{+w}{ }\PY{o}{=}\PY{+w}{ }\PY{p}{[}\PY{l+s+s2}{\PYZdq{}}\PY{l+s+s2}{https://prefix.dev/workshop\PYZhy{}pichi\PYZhy{}\PYZlt{}nombre\PYZgt{}}\PY{l+s+s2}{\PYZdq{}}\PY{p}{,}
\PY{+w}{            }\PY{l+s+s2}{\PYZdq{}}\PY{l+s+s2}{robostack\PYZhy{}jazzy}\PY{l+s+s2}{\PYZdq{}}\PY{p}{,}\PY{+w}{ }\PY{l+s+s2}{\PYZdq{}}\PY{l+s+s2}{conda\PYZhy{}forge}\PY{l+s+s2}{\PYZdq{}}\PY{p}{]}
\PY{n}{platforms}\PY{+w}{ }\PY{o}{=}\PY{+w}{ }\PY{p}{[}\PY{l+s+s2}{\PYZdq{}}\PY{l+s+s2}{linux\PYZhy{}64}\PY{l+s+s2}{\PYZdq{}}\PY{p}{]}
\PY{k}{[}\PY{k}{dependencies}\PY{k}{]}
\PY{n}{ros\PYZhy{}jazzy\PYZhy{}ros\PYZhy{}base}\PY{+w}{ }\PY{o}{=}\PY{+w}{ }\PY{l+s+s2}{\PYZdq{}}\PY{l+s+s2}{*}\PY{l+s+s2}{\PYZdq{}}
\PY{n}{colcon\PYZhy{}common\PYZhy{}extensions}\PY{+w}{ }\PY{o}{=}\PY{+w}{ }\PY{l+s+s2}{\PYZdq{}}\PY{l+s+s2}{*}\PY{l+s+s2}{\PYZdq{}}
\PY{n}{hola\PYZhy{}pichi}\PY{+w}{ }\PY{o}{=}\PY{+w}{ }\PY{l+s+s2}{\PYZdq{}}\PY{l+s+s2}{*}\PY{l+s+s2}{\PYZdq{}}
\PY{k}{[}\PY{k}{tasks}\PY{k}{]}
\PY{n}{build}\PY{+w}{ }\PY{o}{=}\PY{+w}{ }\PY{l+s+s2}{\PYZdq{}}\PY{l+s+s2}{colcon build \PYZhy{}\PYZhy{}symlink\PYZhy{}install}\PY{l+s+s2}{\PYZdq{}}
\PY{k}{[}\PY{k}{activation}\PY{k}{]}
\PY{n}{scripts}\PY{+w}{ }\PY{o}{=}\PY{+w}{ }\PY{p}{[}\PY{l+s+s2}{\PYZdq{}}\PY{l+s+s2}{install/setup.sh}\PY{l+s+s2}{\PYZdq{}}\PY{p}{]}
\end{Verbatim}
\end{filelisting}
\begin{terminal}
pixi install && pixi shell
hola-pichi
\end{terminal}
\end{frame}

\begin{frame}
\begin{center}
\resizebox{\linewidth}{!}{%
\begin{tikzpicture}[node distance=0.9cm, every node/.style={font=\footnotesize}]
\node[draw, rounded corners, fill=green!15, minimum width=2.6cm, minimum height=0.8cm] (a) {receta};
\node[draw, rounded corners, fill=green!20, minimum width=2.6cm, minimum height=0.8cm, right=0.9cm of a] (b) {rattler-build};
\node[draw, rounded corners, fill=orange!15, minimum width=2.6cm, minimum height=0.8cm, right=0.9cm of b] (c) {prefix.dev};
\node[draw, rounded corners, fill=blue!12, minimum width=2.6cm, minimum height=0.8cm, right=0.9cm of c] (d) {\lstinline|pixi.toml|};
\node[draw, rounded corners, fill=blue!20, minimum width=2.6cm, minimum height=0.8cm, right=0.9cm of d] (e) {workspace ROS~2};
\draw[-Latex, thick] (a) -- (b);
\draw[-Latex, thick] (b) -- (c);
\draw[-Latex, thick] (c) -- (d);
\draw[-Latex, thick] (d) -- (e);
\end{tikzpicture}%
}
\end{center}
\vfill
\begin{center}
\Large Habéis recorrido el mismo camino que una build farm real.
\end{center}
\vfill
\normalsize Receta $\rightarrow$ rattler-build $\rightarrow$ prefix.dev $\rightarrow$ \lstinline|pixi.toml| $\rightarrow$ mezclado con ROS~2 oficial de RoboStack. Todo con lockfile, sin \texttt{sudo}, sin Docker, multiplataforma.
\end{frame}

\begin{frame}{Y así lo hacemos de verdad, en producción}
\begin{center}
\normalsize \texttt{IntelligentRoboticsLabs}, para distribuir\\ PlanSys2 y EasyNav como paquetes ROS reproducibles.
\end{center}
\vfill
\begin{block}{El patrón: forkear, no reinventar}
\texttt{ros-rolling} y \texttt{ros-jazzy} son \textbf{forks directos} de \texttt{RoboStack/ros-rolling} y \texttt{RoboStack/ros-jazzy}.
\end{block}
\vfill
Reutilizan el 100\% del motor (Pixi + Vinca + rattler-build + CI + patches) que acabáis de usar en pequeño. Lo único que añaden: \textbf{un fichero que dice qué construir}.
\end{frame}

\begin{frame}[fragile]{\texttt{plansys2\_subset/vinca.yaml}}
\begin{filelisting}
\begin{Verbatim}[commandchars=\\\{\}]
\PY{n+nt}{ros\PYZus{}distro}\PY{p}{:}\PY{+w}{ }\PY{l+lScalar+lScalarPlain}{rolling}
\PY{n+nt}{conda\PYZus{}index}\PY{p}{:}
\PY{+w}{  }\PY{p+pIndicator}{\PYZhy{}}\PY{+w}{ }\PY{l+lScalar+lScalarPlain}{robostack.yaml}
\PY{+w}{  }\PY{p+pIndicator}{\PYZhy{}}\PY{+w}{ }\PY{l+lScalar+lScalarPlain}{packages\PYZhy{}ignore.yaml}

\PY{c+c1}{\PYZsh{} Los paquetes ya en robostack\PYZhy{}rolling se dan por existentes}
\PY{n+nt}{skip\PYZus{}existing}\PY{p}{:}
\PY{+w}{  }\PY{p+pIndicator}{\PYZhy{}}\PY{+w}{ }\PY{l+lScalar+lScalarPlain}{https://prefix.dev/robostack\PYZhy{}rolling/}

\PY{c+c1}{\PYZsh{} Solo estos paquetes (sus dependencias se resuelven solas)}
\PY{n+nt}{packages\PYZus{}select\PYZus{}by\PYZus{}deps}\PY{p}{:}
\PY{+w}{  }\PY{p+pIndicator}{\PYZhy{}}\PY{+w}{ }\PY{l+lScalar+lScalarPlain}{plansys2\PYZus{}bringup}
\PY{+w}{  }\PY{p+pIndicator}{\PYZhy{}}\PY{+w}{ }\PY{l+lScalar+lScalarPlain}{plansys2\PYZus{}core}
\PY{+w}{  }\PY{p+pIndicator}{\PYZhy{}}\PY{+w}{ }\PY{l+lScalar+lScalarPlain}{plansys2\PYZus{}domain\PYZus{}expert}
\PY{+w}{  }\PY{p+pIndicator}{\PYZhy{}}\PY{+w}{ }\PY{l+lScalar+lScalarPlain}{plansys2\PYZus{}executor}
\PY{+w}{  }\PY{c+c1}{\PYZsh{} ... resto de paquetes plansys2\PYZus{}*}
\PY{+w}{  }\PY{p+pIndicator}{\PYZhy{}}\PY{+w}{ }\PY{l+lScalar+lScalarPlain}{popf}
\PY{+w}{  }\PY{p+pIndicator}{\PYZhy{}}\PY{+w}{ }\PY{l+lScalar+lScalarPlain}{rclcpp\PYZus{}cascade\PYZus{}lifecycle}
\end{Verbatim}
\end{filelisting}
\end{frame}

\begin{frame}[fragile]{El atajo, y el flujo manual detrás}
\begin{terminal}
cd ros-rolling
pixi install
pixi run plansys2-all
\end{terminal}
\vfill
\small Por debajo (idéntico al runbook genérico, apuntando al subset con \texttt{-d}):
\begin{terminal}
pixi run vinca -d ./plansys2_subset --platform linux-64 -m
pixi run rattler-build build --recipe-dir ./recipes \
  -c https://prefix.dev/robostack-rolling --skip-existing
pixi run rattler-index fs ./output --force
pixi run rattler-build upload prefix --channel irl-rolling ...
\end{terminal}
\end{frame}

\begin{frame}[fragile]{Cómo lo consume alguien}
\begin{filelisting}
\begin{Verbatim}[commandchars=\\\{\}]
\PY{k}{[}\PY{k}{workspace}\PY{k}{]}
\PY{n}{name}\PY{+w}{ }\PY{o}{=}\PY{+w}{ }\PY{l+s+s2}{\PYZdq{}}\PY{l+s+s2}{plansys2\PYZus{}ws}\PY{l+s+s2}{\PYZdq{}}
\PY{n}{channels}\PY{+w}{ }\PY{o}{=}\PY{+w}{ }\PY{p}{[}
\PY{+w}{  }\PY{l+s+s2}{\PYZdq{}}\PY{l+s+s2}{file:///.../ros\PYZhy{}rolling/output}\PY{l+s+s2}{\PYZdq{}}\PY{p}{,}\PY{+w}{       }\PY{c+c1}{\PYZsh{} local, para iterar}
\PY{+w}{  }\PY{l+s+s2}{\PYZdq{}}\PY{l+s+s2}{https://prefix.dev/irl\PYZhy{}rolling}\PY{l+s+s2}{\PYZdq{}}\PY{p}{,}\PY{+w}{        }\PY{c+c1}{\PYZsh{} canal propio}
\PY{+w}{  }\PY{l+s+s2}{\PYZdq{}}\PY{l+s+s2}{https://prefix.dev/robostack\PYZhy{}rolling}\PY{l+s+s2}{\PYZdq{}}\PY{p}{,}\PY{+w}{  }\PY{c+c1}{\PYZsh{} ROS 2 oficial}
\PY{+w}{  }\PY{l+s+s2}{\PYZdq{}}\PY{l+s+s2}{https://prefix.dev/conda\PYZhy{}forge}\PY{l+s+s2}{\PYZdq{}}\PY{p}{,}
\PY{p}{]}
\PY{n}{platforms}\PY{+w}{ }\PY{o}{=}\PY{+w}{ }\PY{p}{[}\PY{l+s+s2}{\PYZdq{}}\PY{l+s+s2}{linux\PYZhy{}64}\PY{l+s+s2}{\PYZdq{}}\PY{p}{]}

\PY{k}{[}\PY{k}{dependencies}\PY{k}{]}
\PY{n}{ros\PYZhy{}rolling\PYZhy{}ros\PYZhy{}base}\PY{+w}{ }\PY{o}{=}\PY{+w}{ }\PY{l+s+s2}{\PYZdq{}}\PY{l+s+s2}{*}\PY{l+s+s2}{\PYZdq{}}
\PY{n}{ros\PYZhy{}rolling\PYZhy{}plansys2\PYZhy{}bringup}\PY{+w}{ }\PY{o}{=}\PY{+w}{ }\PY{l+s+s2}{\PYZdq{}}\PY{l+s+s2}{*}\PY{l+s+s2}{\PYZdq{}}
\PY{n}{ros\PYZhy{}rolling\PYZhy{}plansys2\PYZhy{}core}\PY{+w}{ }\PY{o}{=}\PY{+w}{ }\PY{l+s+s2}{\PYZdq{}}\PY{l+s+s2}{*}\PY{l+s+s2}{\PYZdq{}}
\PY{n}{ros\PYZhy{}rolling\PYZhy{}popf}\PY{+w}{ }\PY{o}{=}\PY{+w}{ }\PY{l+s+s2}{\PYZdq{}}\PY{l+s+s2}{*}\PY{l+s+s2}{\PYZdq{}}
\end{Verbatim}
\end{filelisting}
\vfill
\small El orden importa: local $\rightarrow$ canal propio $\rightarrow$ RoboStack oficial $\rightarrow$ conda-forge.
\end{frame}

\begin{frame}{Esto funciona de verdad}
\begin{itemize}
    \item PlanSys2 + simulación de Nav2, 4 nodos coordinados (\texttt{rmw\_zenohd}, \texttt{nav2\_sim\_node}, launch con árbol de comportamiento, controlador).
    \item Reproducible con \lstinline|pixi install| + \lstinline|pixi run build|.
\end{itemize}
\vfill
\begin{alertblock}{Aviso honesto}
La publicación a prefix.dev es \textbf{manual/local} — no está automatizada en ningún workflow de CI comprometido en el repo.
\end{alertblock}
\end{frame}

\begin{frame}{Para empresas e instituciones}
\begin{center}
\large Si tenéis software propio que queréis distribuir\\ como paquetes ROS reproducibles:
\end{center}
\vfill
\begin{enumerate}
    \item Fork del repo de RoboStack de vuestra distro.
    \item Vuestro propio \lstinline|<algo>_subset/vinca.yaml| seleccionando qué construir.
    \item Reutilizad el resto de la infraestructura tal cual.
\end{enumerate}
\vfill
\centering \normalsize No hace falta reinventar nada de esto.
\end{frame}

\begin{frame}{Reto final (si sobra tiempo)}
\begin{itemize}
    \item Tomad \lstinline|mi_paquete| del bloque 5 (el nodo ROS~2 propio).
    \item Convertidlo en una receta \texttt{rattler-build} (dependiendo de \lstinline|ros-jazzy-rclpy|).
    \item Publicadlo en vuestro mismo canal.
\end{itemize}
\vfill
\centering Cerrando el círculo completo con un paquete ROS real.
\end{frame}

%=================================================
\section[Cierre]{Cierre}
%=================================================

\begin{frame}{Resumen de timing}
\centering
\begin{tabular}{lcc}
\toprule
\textbf{Bloque} & \textbf{Duración} & \textbf{Tipo} \\
\midrule
1. Introducción & 10 min & Teoría \\
2. ¿Qué es Pixi? & 30 min & Teoría \\
3. Instalación y uso básico de Pixi & 30 min & Práctica \\
4. ¿Qué es RoboStack? & 15 min & Teoría \\
5. RoboStack + colcon & 20 min$^*$ & Práctica \\
6. Infraestructura de RoboStack & 15 min & Teoría \\
7. Tu propia Buildfarm & 35 min & Práctica \\
\midrule
\textbf{Núcleo} & \textbf{155 min} & \\
\textbf{Total esperado} & \textbf{\textasciitilde 165-170 min (\textasciitilde 3h)} & \\
\bottomrule
\end{tabular}
\vfill
\footnotesize $^*$ + 10-12 min del ejercicio Nav2 + \texttt{pixi-ros} (demo del instructor, build lanzado de antemano). Margen restante hasta las 3h: ejercicio de varias distros ROS~2 y reto final del bloque 7, ambos opcionales, más preguntas.
\end{frame}

\begin{frame}{Recursos}
\small
\begin{itemize}
    \item Pixi: \url{https://pixi.sh/} — \url{https://prefix.dev/}
    \item RoboStack: \url{https://robostack.github.io/} — \url{https://github.com/RoboStack}
    \item pixi-ros: \url{https://github.com/ruben-arts/pixi-ros} — Nav2: \url{https://github.com/ros-navigation/navigation2}
    \item Buildfarm propia: \url{https://github.com/IntelligentRoboticsLabs/ros-rolling} — \url{https://github.com/IntelligentRoboticsLabs/ros-jazzy}
    \item Paper original: Fischer et al., \emph{``A RoboStack Tutorial''}, IEEE RAM 2021 (arXiv:2104.12910)
\end{itemize}
\end{frame}

\begin{frame}
  \centering
  \Large Preguntas
\end{frame}

\end{document}
